\subsection*{Agent Message Format (JSON Schema)}
\phantomsection
\addcontentsline{toc}{subsection}{Agent Message Format (JSON Schema)}

\noindent
This appendix summarises the structured JSON messages exchanged between
the LLM agents and the LangGraph orchestration layer.
The environment's natural-language state is emitted in Chinese; field meanings
below are given in English for clarity.

\subsubsection*{Environment state description (natural language input)}

The environment builds a single string passed to every LLM call as
\texttt{\{state\_description\}}. Each line has the following semantics
(runtime text is Chinese; structure is fixed):

\begin{itemize}\setlength{\itemsep}{2pt}
  \item Header: current decision step index and total system bandwidth (MHz).
  \item One block per slice \texttt{eMBB}, \texttt{URLLC}, \texttt{mMTC}:
        user count, average CQI, buffer occupancy in $[0,1]$, current bandwidth
        fraction and absolute MHz, sliding-window SLA satisfaction rate in
        $[0,1]$.
\end{itemize}

\subsubsection*{Slice manager agent --- proposal object}

Parsed by \texttt{JsonOutputParser}; keys are normalised before downstream use.
Intended shape:

\begin{lstlisting}[basicstyle=\footnotesize\ttfamily, breaklines=true, columns=fullflexible, frame=single]
{
  "slice": "eMBB" | "URLLC" | "mMTC",   // injected after parsing
  "requested": <float>,                // requested bandwidth fraction
  "minimum": <float>,                  // >= 0.05 after validation
  "priority": "high" | "medium" | "low",
  "justification": "<string>"
}
\end{lstlisting}

\noindent
The \texttt{slice} field is appended in code so downstream components always
know the source slice. Defaults on parse failure are defined per slice.

\subsubsection*{Global coordinator agent --- evaluation object}

Returned by the coordinator LLM; allocation is renormalised to sum to~1.

\begin{lstlisting}[basicstyle=\footnotesize\ttfamily, breaklines=true, columns=fullflexible, frame=single]
{
  "compatible": <bool>,
  "allocation": {
    "eMBB": <float>,
    "URLLC": <float>,
    "mMTC": <float>
  },
  "feedback": "<string>"
}
\end{lstlisting}

\noindent
When building the coordinator prompt, each proposal is flattened to one line of
text for \texttt{\{proposals\_text\}}, e.g.\
\texttt{[eMBB] requested=0.xx, minimum=0.xx, priority=..., justification=...}.

\subsubsection*{Single-agent baseline (M4) --- allocation object}

Direct mapping from JSON to the three-dimensional action (then clipped to
$\geq 0.05$ and normalised):

\begin{lstlisting}[basicstyle=\footnotesize\ttfamily, breaklines=true, columns=fullflexible, frame=single]
{ "eMBB": <float>, "URLLC": <float>, "mMTC": <float> }
\end{lstlisting}

\subsubsection*{LangGraph multi-agent state (\texttt{MASState})}

Typed dictionary carried across nodes in the LangGraph workflow:

\begin{table}[h]
\centering
\small
\begin{tabular}{@{}p{3.6cm}p{9.8cm}@{}}
\hline
\textbf{Field} & \textbf{Role} \\
\hline
\texttt{state\_description} &
Natural-language snapshot from the environment (see above). \\
\texttt{proposals} &
List of slice proposal dicts (latest round). \\
\texttt{coordinator\_result} &
Coordinator evaluation dict; \texttt{compatible} drives routing. \\
\texttt{negotiation\_round} &
Integer; incremented after each negotiation node visit. \\
\texttt{max\_rounds} &
Cap on negotiation iterations (default 3). \\
\texttt{strategy} &
\texttt{"priority"} or \texttt{"proportional"}: selects
Priority Arbitration vs.\ Proportional Compromise. \\
\texttt{final\_allocation} &
Final dict \texttt{\{"eMBB": float, "URLLC": float, "mMTC": float\}} passed to the env. \\
\texttt{total\_tokens} &
Heuristic token-cost accumulator for logging. \\
\texttt{is\_resolved} &
Set when negotiation applies a resolution. \\
\hline
\end{tabular}
\end{table}

\subsubsection*{Routing and final action}

If \texttt{coordinator\_result.compatible} is true, the graph goes to
the finalise node and uses the coordinator's \texttt{allocation}. Otherwise the
negotiation node computes an allocation from the current
\texttt{proposals} using the selected strategy; the graph may loop back to slice
agents until \texttt{negotiation\_round >= max\_rounds} or resolution. The
finalise node outputs a dict with keys \texttt{eMBB}, \texttt{URLLC},
\texttt{mMTC} (floats, interpreted as fractions, then clipped and
renormalised to sum to~1).

The no-negotiation variant (M5-no-neg) omits the negotiation loop; the
final action is taken from the coordinator's \texttt{allocation} only.
